\textbf{Lösung zu Übung 3, Aufgabe 1}

\begin{loesung}
a) EA für 7, 5 \qquad $z = 5x + 3 = 7y + 4$
\begin{align*}
7 &= 1 \cdot 5 + 2\\
5 &= 2 \cdot 2 + 1\\
2 &= 2 \cdot 1 + 0
\end{align*}
Der letzte Rest $\neq 0$ ist $r_3 = 1$. Jetzt gibt es zwei Lösingsvarianten.

Variante 1: EA rückwärts
\begin{align*}
\Rightarrow\quad 1 &= 5 - 2 \cdot 2\\
&= 5 - 2 \cdot (7 - 1 \cdot 5)\\
&= 3 \cdot 5 - 2 \cdot 7 = 4 - 3\\
\Rightarrow\quad z &= 3 \cdot 5 + 3 = 2 \cdot 7 + 4 = 18 \pmod{35}
\end{align*}
Variante 2: mittels der Matrix $Q$
\[
\begin{pmatrix} m \\ n \end{pmatrix} = Q \begin{pmatrix} r_3 \\ 0 \end{pmatrix}, \quad
Q = \begin{pmatrix} 1 & 1 \\ 1 & 0 \end{pmatrix}\begin{pmatrix} 2 & 1 \\ 1 & 0 \end{pmatrix}\begin{pmatrix} 2 & 1 \\ 1 & 0 \end{pmatrix}
= \begin{pmatrix} 3 & 1 \\ 2 & 1 \end{pmatrix}\begin{pmatrix} 2 & 1 \\ 1 & 0 \end{pmatrix}
= \begin{pmatrix} 7 & 3 \\ 5 & 2 \end{pmatrix}
\]
\[
\operatorname{Det} Q = 7 \cdot 2 - 5 \cdot 3 = -1 = a - b,
\]
\[
\Rightarrow\quad z = 7 \cdot 2 + 4 = 5 \cdot 3 + 3 = 18 \pmod{35}
\]

b) EA für 12, 5
\begin{align*}
12 &= 2 \cdot 5 + 2\\
5 &= 2 \cdot 2 + 1
\end{align*}
EA rückwärts \qquad $b - a = 2 - (-1) = 3$
\begin{align*}
\Rightarrow\quad 1 &= 5 - 2 \cdot 2 = 5 - 2 \cdot (12 - 2 \cdot 5)\\
&= 5 \cdot 5 - 2 \cdot 12\\
2 - (-1) &= 5 \cdot 5 \cdot 3 - 2 \cdot 12 \cdot 3\\
\Rightarrow\quad z &= 2 - 5 \cdot 5 \cdot 3 = (-1) - 2 \cdot 12 \cdot 3 = -73 \pmod{60}
\end{align*}
\end{loesung}
